% Рисунок из тетради: |x| сглаживается в окрестности нуля, производная -> sign x.
\begin{figure}[htbp]
\centering
\begin{tikzpicture}
\begin{groupplot}[
  group style={group size=2 by 1, horizontal sep=1.8cm},
  width=7.2cm, height=4.6cm,
  axis lines=middle,
  samples=120, clip=false,
  ytick=\empty,
]
% --- слева: f = |x| и сглаженная f_k
\nextgroupplot[xmin=-1.3, xmax=1.3, ymin=-0.62, ymax=1.45,
               xtick={-0.3,0.3}, xticklabels={$-\tfrac1k$,$\tfrac1k$}]
  \addplot[very thick, blue!60!black, domain=-1.2:0] {-x};
  \addplot[very thick, blue!60!black, domain=0:1.2] {x};
  % сглаженная ветвь: парабола, склеенная с |x| в точках +-1/k
  \addplot[very thick, red!70!black, domain=-0.3:0.3, samples=60]
    {(x*x + 0.09)/(2*0.3)};
  \draw[dashed, gray] (axis cs:-0.3,0) -- (axis cs:-0.3,0.3);
  \draw[dashed, gray] (axis cs: 0.3,0) -- (axis cs: 0.3,0.3);
  \node[blue!60!black, anchor=south east] at (axis cs:-0.75,0.8) {$f=|x|$};
  \node[red!70!black, anchor=west] at (axis cs:0.33,0.18) {\footnotesize $f_k$};
  \node[anchor=north] at (axis cs:0,-0.42) {\footnotesize $f_k \to f$};

% --- справа: производные
\nextgroupplot[xmin=-1.3, xmax=1.3, ymin=-1.95, ymax=1.6,
               xtick={-0.3,0.3}, xticklabels={$-\tfrac1k$,$\tfrac1k$}]
  \addplot[very thick, blue!60!black] coordinates {(-1.2,-1) (-0.3,-1)};
  \addplot[very thick, blue!60!black] coordinates {(0.3,1) (1.2,1)};
  \addplot[very thick, red!70!black] coordinates {(-0.3,-1) (0.3,1)};
  \draw[dashed, gray] (axis cs:-1.3,1) -- (axis cs:1.3,1);
  \draw[dashed, gray] (axis cs:-1.3,-1) -- (axis cs:1.3,-1);
  \node[red!70!black, anchor=west] at (axis cs:0.35,-0.55) {\footnotesize $f_k'$};
  \node[anchor=north] at (axis cs:0,-1.35) {\footnotesize $f_k' \to \sgn x$};
\end{groupplot}
\end{tikzpicture}
\caption{Пример к признаку: $|x|$ сглаживается на $[-\tfrac1k,\tfrac1k]$, так что
$f_k \to f$ и $f_k' \to \sgn x$ в $\Lloc$. По признаку существования производной
предельная функция $\sgn x$ и есть обобщённая производная: $(|x|)' = \sgn x$.}
\label{fig:abs-x}
\end{figure}
